% Balambika specimen and reference setup for XeLaTeX.
% Build: xelatex xelatex-specimen.tex
%
% XeTeX shapes each font instance with a single OpenType script, so
% Devanagari must be set with a font instance that uses Script=Devanagari,
% even though Balambika contains all scripts in one file. Polyglossia's
% \textsanskrit / sanskrit environment does the switching and also selects
% Sanskrit hyphenation.
\documentclass[11pt]{article}
\usepackage[a4paper, margin=25mm]{geometry}
\usepackage{fontspec}
\usepackage{polyglossia}
\setdefaultlanguage{bulgarian}
\setotherlanguages{english, sanskrit}

\defaultfontfeatures{Path=../fonts/ttf/, Extension=.ttf}

% Latin, IAST and Cyrillic.
\setmainfont{Balambika}[
  UprightFont    = *-Regular,
  ItalicFont     = *-Italic,
  BoldFont       = *-Bold,
  BoldItalicFont = *-BoldItalic,
  FontFace       = {mb}{n}{*-Medium},
  FontFace       = {mb}{it}{*-MediumItalic},
]

% Devanagari. XeTeX measures the word space from a lone U+0020, so the
% font's contextual Devanagari space (calt) never applies here; WordSpace
% restores Shobhika's spacing: 450/259 (Regular), 500/259 (Medium, Bold).
\newfontfamily\devanagarifont{Balambika}[
  Script         = Devanagari,
  WordSpace      = 1.737,
  BoldFeatures   = {WordSpace = 1.931},
  UprightFont    = *-Regular,
  ItalicFont     = *-Italic,
  BoldFont       = *-Bold,
  BoldItalicFont = *-BoldItalic,
  FontFace       = {mb}{n}{Font = *-Medium, WordSpace = 1.931},
  FontFace       = {mb}{it}{Font = *-MediumItalic, WordSpace = 1.931},
]
\setsansfont{Balambika}[UprightFont=*-Regular, BoldFont=*-Bold, ItalicFont=*-Italic, BoldItalicFont=*-BoldItalic]

% Explicit Latin-script instance for IAST: inside Sanskrit, polyglossia
% would otherwise select the Devanagari instance even for Latin script.
\newfontfamily\iastfont{Balambika}[
  UprightFont    = *-Regular,
  ItalicFont     = *-Italic,
  BoldFont       = *-Bold,
  BoldItalicFont = *-BoldItalic,
]

% Shorthands: Devanagari word, IAST transliteration.
\newcommand\dev[1]{\textsanskrit{#1}}
\newcommand\iast[1]{\textlang{sa-Latn}{\iastfont\itshape #1}}

\linespread{1.15}
\setlength\parindent{0pt}
\setlength\parskip{6pt}

\begin{document}

{\Huge Balambika}\par
{\small Shobhika + Common Serif (108\%) · SIL Open Font License 1.1}

\section*{Ṛgveda 1.1.1}

\begin{sanskrit}
अ॒ग्निमी॑ळे पु॒रोहि॑तं य॒ज्ञस्य॑ दे॒वमृ॒त्विज॑म् । होता॑रं रत्न॒धात॑मम् ॥
\end{sanskrit}

\iast{agním īḷe puróhitaṃ yajñásya devám ṛtvíjam | hótāraṃ ratnadhā́tamam ||}

Агни възхвалявам — поставения отпред жрец, бога на жертвата,
свещенослужителя, призоваващия, най-щедрия даряващ съкровища.

\section*{Mixed scripts}

\begin{english}
The term \dev{अग्निः} (\iast{agniḥ}) denotes both the sacrificial fire and the
deity; cf. \dev{यज्ञः} (\iast{yajñaḥ}) and \dev{होता} (\iast{hotā}).
\end{english}

Терминът \dev{अग्निः} (\iast{agniḥ}) означава едновременно жертвения огън и
божеството; срв. \dev{यज्ञः} (\iast{yajñaḥ}) и \dev{होता} (\iast{hotā}).

\section*{Bold Devanagari in regular text}

\begin{english}
The term \textbf{\dev{अग्निः}} (\iast{agniḥ}) denotes both the sacrificial fire
and the deity; cf. \textbf{\dev{यज्ञः}} (\iast{yajñaḥ}) and
\textbf{\dev{होता}} (\iast{hotā}).
\end{english}

Терминът \textbf{\dev{अग्निः}} (\iast{agniḥ}) означава едновременно жертвения
огън и божеството; срв. \textbf{\dev{यज्ञः}} (\iast{yajñaḥ}) и
\textbf{\dev{होता}} (\iast{hotā}).

\section*{Weights}

\begin{tabular}{@{}ll@{}}
Regular & \dev{श्रीमाता श्रीमहाराज्ञी} — \textit{śrīmātā śrīmahārājñī} — Шри Видя \\
Medium  & {\fontseries{mb}\selectfont \dev{श्रीमाता श्रीमहाराज्ञी} — \textit{śrīmātā} — Шри Видя} \\
Bold    & \textbf{\dev{श्रीमाता श्रीमहाराज्ञी} — \textit{śrīmātā} — Шри Видя} \\
\end{tabular}

\section*{Coverage}

IAST: \iast{ā ī ū ṛ ṝ ḷ ḹ ṃ ṁ ḥ ṅ ñ ṭ ḍ ṇ ś ṣ m̐}\\
{\upshape ĀĪŪṚṜḶḸṂṀḤṄÑṬḌṆŚṢ — ā́ ī́ ṛ́ — ṚGVEDA SAṂHITĀ}

Кирилица: вдигнахме жълтия гълъб, щом дойде юли; дълг, гъба, ѝ.
\textit{вдигнахме жълтия гълъб, щом дойде юли; дълг, гъба, ѝ.}

\begin{sanskrit}
क्ष्म्ल्व्र्यूं ह्रीं श्रीं क्लीं ऐं सौः · र्क र्कि द्ध द्व ह्म ह्य ङ्क्ष ष्ट्र त्त्व · ०१२३४५६७८९ ·
ꣳ ꣴ ᳵ ᳶ ᳐ ᳑ ᳒ ᳚ ᳛ ँ ं ः ऽ ॐ
\end{sanskrit}

\end{document}
